\documentclass[11pt,UTF8,fontset=fandol]{ctexart}
\usepackage{amsmath,amssymb,amsfonts}
\usepackage{graphicx,float,caption,subcaption}
\usepackage{booktabs,array,multirow,makecell}
\usepackage{geometry,enumitem,xcolor}
\usepackage[hidelinks]{hyperref}
\usepackage{url}
\geometry{margin=2.3cm}
\graphicspath{{/Users/yye/project/论文/assets/uitars/}}
\captionsetup{font=small,labelfont=bf}
\setlength{\parskip}{0.4em}

% 知识卡片环境
\definecolor{cardbg}{RGB}{244,247,252}
\definecolor{cardline}{RGB}{120,150,200}
\newenvironment{kcard}[1]{%
  \par\smallskip\noindent
  \begin{tabular}{@{}p{0.985\linewidth}@{}}
  \rule{0pt}{1.1em}\textbf{\textcolor{cardline}{#1}}\\[0.2em]
}{%
  \\[0.3em]\end{tabular}\par\smallskip
}

\title{\textbf{UI-TARS：开创纯截图驱动的原生 GUI 操作智能体}}
\author{Yujia Qin, Yining Ye, Junjie Fang, Haoming Wang, Shihao Liang, 等\\
ByteDance Seed（字节跳动）\,/\,Tsinghua University（清华大学）}
\date{}

\begin{document}
\maketitle

\begin{center}
\fbox{\parbox{0.95\linewidth}{\small
\textbf{英文原标题：}UI-TARS: Pioneering Automated GUI Interaction with Native Agents\\
\textbf{发表：}arXiv 2501.12326（2025 年 1 月，字节跳动 ByteDance Seed）\quad\textbf{引用数：}449（Semantic Scholar，2026-06-07）\\
\textbf{代码：}\url{https://github.com/bytedance/UI-TARS}\\
\textbf{一句话定位：}做一个\textbf{“只看屏幕截图就能像人一样操作电脑/手机”}的端到端原生 GUI 智能体——不读 HTML、不用辅助接口，看界面 $\rightarrow$ 想清楚 $\rightarrow$ 点哪/输什么，全部由\textbf{一个模型}搞定，并能从自己的操作历史里反思改进。
}}
\end{center}

\begin{abstract}
\noindent 本文介绍 UI-TARS，一个\textbf{原生 GUI 智能体模型（native GUI agent model）}，它\textbf{仅以屏幕截图作为输入}，并执行类人交互（如键盘、鼠标操作）。与当前主流的、依赖被层层包装的商业模型（如 GPT-4o）外加专家精心设计的提示词和工作流的智能体框架不同，UI-TARS 是一个\textbf{端到端模型}，性能却超过了这些复杂框架。实验表明其性能领先：UI-TARS 在 10 多个评测感知、定位（grounding）与 GUI 任务执行能力的基准上取得 SOTA。尤其在 OSWorld 基准上，UI-TARS 在 50 步预算下得分 24.6、15 步下得分 22.7，分别超过 Claude 的 22.0 和 14.9；在 AndroidWorld 上得分 46.6，超过 GPT-4o 的 34.5。UI-TARS 包含若干关键创新：（1）\textbf{增强感知}：利用大规模 GUI 截图数据集，实现对 UI 元素的上下文感知理解与精确描述；（2）\textbf{统一动作建模}：把动作标准化到一个跨平台统一的动作空间，并通过大规模动作轨迹实现精确定位与交互；（3）\textbf{System-2 推理}：把审慎推理融入多步决策，涵盖任务分解、反思思考、里程碑识别等多种推理模式；（4）\textbf{基于反思式在线轨迹的迭代训练}：通过在数百台虚拟机上自动采集、筛选并反思式精炼新的交互轨迹来突破数据瓶颈。通过迭代训练与反思微调，UI-TARS 能持续从自身错误中学习，并在极少人工干预下适应未曾见过的情况。我们还分析了 GUI 智能体的演进路径，以指导该领域的进一步发展。UI-TARS 已在 \url{https://github.com/bytedance/UI-TARS} 开源。
\end{abstract}

% =========================================================
\section{小白知识卡：读这篇前要懂的背景}

\begin{kcard}{GUI 智能体（GUI Agent）}
能自动操作图形界面（按钮、输入框、菜单……）来完成任务的 AI，比如让它“帮我订一张明天去北京的机票”，它自己点来点去把任务做完。
\end{kcard}

\begin{kcard}{原生智能体模型 vs.\ 智能体框架（Native Agent Model vs.\ Agent Framework）}
\textbf{框架}=把一个强模型（如 GPT-4o）当大脑，外面再用人手写的提示词、脚本、各种工具模块（看图的、记忆的、规划的）拼起来——\textbf{设计驱动}，环境一变就容易崩。\textbf{原生模型}=把所有能力（看、想、记、做）都\textbf{学进同一个模型的参数里}，端到端，\textbf{数据驱动}，能持续学习适应。本文主张后者是未来。
\end{kcard}

\begin{kcard}{纯视觉 / 截图驱动（Pure-Vision）}
模型\textbf{只看界面截图}（像人用眼睛看屏幕），不去读网页的 HTML 或系统的无障碍树（accessibility tree）。好处：不用系统权限、跨平台通用、更接近人的认知方式。
\end{kcard}

\begin{kcard}{定位 / 接地（Grounding）}
把“我想点的那个东西”精确翻译成屏幕上的\textbf{坐标}（x, y）。比如“点右上角那个写着 Submit 的红按钮”$\rightarrow$ 输出像素坐标。这是截图智能体能不能真正“点对地方”的关键。
\end{kcard}

\begin{kcard}{System-1 与 System-2 思维（快思考与慢思考）}
借自心理学：\textbf{System-1}=快、直觉、不假思索（如熟练地点个常见按钮）；\textbf{System-2}=慢、刻意、一步步分析（如规划整个流程、出错后反思纠错）。本文重点是给模型装上 System-2：动手前先生成一段\textbf{“思考（thought）”}。
\end{kcard}

\begin{kcard}{VLM（视觉语言模型，Vision-Language Model）}
能同时看图又能读写文字的大模型。本文的底座是 \textbf{Qwen-2-VL}（7B 和 72B 两个规模）。
\end{kcard}

\begin{kcard}{行为克隆（Behavior Cloning）}
最朴素的“照着人怎么做来学”：给模型看“专家在某界面做了什么动作”，让它模仿。本文用它来学习多步操作轨迹。
\end{kcard}

\begin{kcard}{SFT / DPO（监督微调 / 直接偏好优化，Direct Preference Optimization）}
\textbf{SFT}=拿“正确示范”让模型照着学。\textbf{DPO}=给模型成对的“好动作 vs.\ 坏动作”，直接告诉它“这个比那个好”，让它更偏向好动作、远离坏动作。本文用 DPO 把“犯错的步骤”也利用起来（SFT 只用对的步骤，浪费了错的）。
\end{kcard}

\begin{kcard}{推理时扩展 / Best-of-N（Inference-time Scaling）}
不改模型参数，靠“多采样几次再挑最好的”来提升表现。本文实验里让模型对同一输入采样 $N=1/16/64$ 个候选，发现 System-2 在 N 大时优势明显。
\end{kcard}

% =========================================================
\section{逐节精读}

\subsection{引言（§1）：为什么要做原生 GUI 智能体}

论文先点出大方向：自主智能体（autonomous agents）要能在极少人工监督下感知环境、决策、执行动作。而让智能体\textbf{无缝操作图形界面（GUI）}是一个关键前沿——因为人类的数字工作充满了按钮、文本框、图片。

\textbf{过去的两大老路子及其毛病：}
\begin{itemize}[leftmargin=1.5em]
\item \textbf{依赖文本表示}（读 HTML、DOM、无障碍树）：平台之间不一致、太啰嗦、可扩展性差，还经常需要系统级权限才能拿到这些信息，限制了通用性。
\item \textbf{模块化的智能体框架}：把功能拆成多个组件（用 GPT-4o 做理解推理、另用工具做定位或记忆），靠人手工拼装。开发某类任务快，但依赖专家知识与任务定制，\textbf{换个陌生任务或环境一变就容易失败}。
\end{itemize}

\textbf{两个关键转向（这就是“原生”的含义）：}（1）从依赖文本 $\rightarrow$ \textbf{纯视觉}（只看截图）；（2）从模块化框架 $\rightarrow$ \textbf{端到端模型}。论文有一句很点题的对比（关键句翻译）：

\begin{quote}\small
“框架是\textbf{设计驱动的}，需要大量人工工程和预定义工作流来维持稳定、防止意外；而智能体模型本质是\textbf{数据驱动的}，能够通过大规模数据和迭代反馈来学习和适应。”
\end{quote}

\textbf{但原生智能体目前为何还不够好？}两大根源：（1）GUI 领域本身难——感知上要理解高信息密度且不断变化的界面；推理规划要能导航和操纵界面；还要用到记忆（参考过去交互）；更要执行\textbf{精确的底层动作}（输出准确的点击坐标、往正确的输入框打字）。（2）从框架转向模型带来\textbf{根本性的数据瓶颈}：模块化框架可以用各组件单独的数据集（好凑），但端到端模型需要\textbf{把感知、推理、记忆、动作整合在一个统一工作流里的数据}，而这种富含人类专家工作流知识的数据\textbf{历史上极其稀缺}。

\textbf{本文的四大核心贡献}（对应后文 §4.2--4.5）：
\begin{enumerate}[leftmargin=1.6em]
\item \textbf{增强 GUI 截图感知}（§4.2）：用专门的解析工具采集网站/应用/系统的截图与元数据，构建大规模数据集，设计五类任务（元素描述、密集描述、状态转移描述、问答、Set-of-Mark 标注提示）。
\item \textbf{统一动作建模与多步执行}（§4.3）：设计跨平台统一动作空间；用标注轨迹 + 标准化的开源数据训练多步执行；用“元素描述配坐标”的大数据提升定位能力。
\item \textbf{面向审慎决策的 System-2 推理}（§4.4）：爬取并精炼约 \textbf{600 万}条 GUI 教程注入推理知识；给所有动作轨迹\textbf{增补“思考（thought）”}，注入任务分解、长期一致性、里程碑识别、试错、反思等推理模式——\textbf{每个动作前先生成显式思考}。
\item \textbf{从经验学习的迭代精炼}（§4.5）：在数百台虚拟机上在线探索、采集、筛选、反思式精炼新轨迹回灌模型，持续迭代；并通过\textbf{反思微调}和 \textbf{Agent DPO} 让模型学会识别并从错误中恢复。
\end{enumerate}

底座是 \textbf{Qwen-2-VL}（7B 与 72B），在约 \textbf{500 亿（50B）token} 上继续训练，得到 UI-TARS-7B 与 UI-TARS-72B。（注：模型名 TARS 致敬电影《星际穿越》里那个机智、善于解决问题的机器人 TARS。）

\subsection{GUI 智能体的演进路径（§2）}

论文把 GUI 智能体的发展按\textbf{“需要多少人工干预、泛化能力多强”}分为四个阶段（见图~\ref{fig:arch} 的脉络，原文为 Figure 2/3）：

\begin{itemize}[leftmargin=1.5em]
\item \textbf{阶段 1：基于规则的智能体（Rule-based）}。如 RPA（机器人流程自动化）系统：把用户指令匹配到预定义规则、调用 API。做固定重复任务有效，但完全靠人定规则、需要直接的 API/系统权限，遇到新情况就抓瞎，无法从环境或经验中学习。
\item \textbf{阶段 2：智能体框架（Agent Framework）}。用 GPT-4 / GPT-4o 这类基础模型当大脑（AutoGPT、LangChain、ReAct、SeeAct 等），靠提示工程与多步推理变灵活。但\textbf{“工作流知识”仍是人手工编码}的（提示词、外部脚本、工具使用启发式），带来三大弊端：\textbf{脆弱且维护成本高}（任务/界面一变就得重写规则）、\textbf{学习范式割裂}（很少用新经验去更新底层模型参数）、\textbf{模块不兼容}（多个模块靠提示或胶水代码协调，一处出错整条链崩，还得专家来调试）。
\item \textbf{阶段 3：原生智能体模型（Native Agent Model）}——\textbf{本文所处的阶段}。把工作流知识\textbf{直接学进模型参数}，端到端统一感知/推理/记忆/动作，数据驱动。优势：整体学习与适应、减少人工工程、统一参数带来强泛化、天然支持在线/终身学习。代表工作：Claude Computer-Use、Aguvis、ShowUI、OS-Atlas 等。
\item \textbf{阶段 4：主动与终身学习智能体（Active \& Lifelong，展望）}。智能体\textbf{自己提出任务、执行、评估结果、自我打分奖励}，靠自驱探索发现新知识、持续提升，几乎不依赖人工标注。与阶段 3 的本质区别在于\textbf{学习过程的自主性}：原生模型仍依赖人，主动智能体自己驱动学习。本文目标是建一个可扩展、数据驱动的原生模型，为这一阶段铺路。
\end{itemize}

\subsection{原生智能体模型的核心能力（§3）}

论文把框架时代的模块“内化”为模型的四大核心能力：

\textbf{感知（Perception）。} 按输入特征把已有工作分三类：\textbf{结构化文本}（HTML/DOM/无障碍树，早期 LLM 智能体被迫这么做）、\textbf{视觉截图}（用 SoM 提示、OCR、GUI 元素检测如 DINO 来理解界面）、\textbf{综合界面建模}（同时用文本、视觉、语义轮廓）。还强调\textbf{实时感知}的重要性（界面是动态的，比如出现加载转圈要识别为“正在处理”）。结论：原生模型应像人一样\textbf{直接通过视觉感知}并把动作精确接地到原始截图上。

\textbf{动作（Action）。} 要点：\textbf{统一且多样的动作空间}（把各平台的 click/type/scroll/drag 抽象成一套通用操作，区分原子动作与组合动作）；\textbf{坐标定位的挑战}（不同设备宽高比、动态内容让精确定位很难）。

\textbf{推理（System-1 \& System-2）。} System-1 快而直觉但只能反应式一步决策；System-2 显式生成中间思考（用 CoT 或 ReAct），包含五个要素：\textbf{任务分解、长期一致性、里程碑识别、试错、反思}。UI-TARS 重点强化 System-2。

\textbf{记忆（Memory）。} 分\textbf{短期记忆}（当前任务的动作历史、当前状态、执行轨迹）与\textbf{长期记忆}（过往任务的执行路径、背景知识）。原生模型把长期经验\textbf{编码进自身参数}，用 In-Context Learning 或 CoT 来激活。

\textbf{能力评测分类（§3.2）}：感知评测、定位评测，以及智能体能力评测——又分\textbf{离线}（静态预定义环境，评单步，用 step-level 指标如 Action-Matching Score）与\textbf{在线}（可交互的动态环境，评能否真正完成整个任务，用 task-level 成功率）。

\subsection{UI-TARS 总体架构（§4.1）}

这是方法部分的总入口，对应图~\ref{fig:arch}（原文 Figure 4）。给定一个初始任务指令，UI-TARS 不断地从设备收到\textbf{观测（observation，即截图）}、执行对应\textbf{动作（action）}。整个过程可形式化为一个交替序列：
\[
(\text{instruction},\,(o_1,a_1),\,(o_2,a_2),\,\cdots,\,(o_n,a_n)),
\]
其中 $o_i$ 是第 $i$ 步的观测（设备截图），$a_i$ 是智能体执行的动作。每一步，模型输入\textbf{任务指令 + 历史交互 $(o_1,a_1,\cdots,o_{i-1},a_{i-1})$ + 当前观测 $o_i$}，输出一个来自预定义动作空间的动作 $a_i$；执行后设备给出下一帧观测，循环往复。

为强化推理、促成更审慎的决策，作者在\textbf{每个动作 $a_i$ 之前插入一段“思考（thought）”$t_i$}，体现 System-2 的反思特性。它是承上启下的关键中间步骤，让智能体在前进前先重新审视之前的动作和观测。这受 ReAct 启发，但比 ReAct 更结构化、更目标导向。于是过程变为：
\[
(\text{instruction},\,(o_1,t_1,a_1),\,(o_2,t_2,a_2),\,\cdots,\,(o_n,t_n,a_n)).
\]

为了在有限 token 预算内（如 32k 序列长度）控制开销，输入只保留\textbf{最近 $N$ 帧观测}，但完整的历史动作与思考作为短期记忆保留。模型迭代预测当前思考 $t_n$ 和动作 $a_n$：
\[
P\big(t_n,a_n \mid \text{instruction},\,t_1,a_1,\cdots,\,(o_{n-i},t_{n-i},a_{n-i})_{i=1}^{N},\,o_n\big).
\]
\textbf{公式拆解：}竖线 $\mid$ 左边是\textbf{要预测的}（这一步的思考+动作），右边是\textbf{已知条件}（任务指令、全部历史思考与动作、最近 $N$ 帧观测、当前帧 $o_n$）。实验中 $N$ 取 5。

\subsection{增强 GUI 感知（§4.2）}

\textbf{两大难点：}（1）\textbf{截图稀缺}——通用自然图像很多，GUI 截图相对少；（2）\textbf{信息密度高、精度要求高}——一张界面常有上百个元素、布局复杂，且很多元素极小（如 $1920\times1080$ 图里 $10\times10$ 像素的图标），难以准确感知和定位。

\textbf{数据采集：}用专门解析工具，一边渲染截图一边自动抽取丰富元数据（元素类型、层级深度、包围框 bounding box、文本内容），结合自动爬取 + 人工辅助探索，覆盖主界面与深层嵌套页面，统一记成 \texttt{(截图, 元素框, 元素元数据)} 结构。采用\textbf{自底向上}的构建思路：先搞懂局部小元素，再整合到整体界面理解。

基于这些数据，设计了\textbf{五类核心任务}（对应原文 Figure 5）：
\begin{enumerate}[leftmargin=1.6em]
\item \textbf{元素描述（Element Description）}：对每个元素给出结构化描述，含元素类型、视觉外观、位置信息、元素功能；训练模型枚举截图里所有可见元素并描述。
\item \textbf{密集描述（Dense Captioning）}：对整张界面做全面细致的描述，不仅是元素本身，还包括空间关系与整体布局（嵌入图片也生成描述），同时尽量减少幻觉。
\item \textbf{状态转移描述（State Transition Captioning）}：识别并描述\textbf{两张连续截图之间的差异}，判断是否发生了某个动作（如鼠标点击、键盘输入），也包含非交互式变化（动画、刷新）。这对细粒度交互理解、动态状态感知很关键。
\item \textbf{问答（QA）}：合成涵盖界面理解、图像解读、元素识别、关系推理的多样问答，提升抽象与推理能力。
\item \textbf{Set-of-Mark（SoM）提示}：按坐标在截图上画出形状/颜色/大小各异的视觉标记，让模型更好地把视觉标记与对应元素关联起来。
\end{enumerate}

\subsection{统一动作建模与定位（§4.3）}

核心方法是\textbf{行为克隆}：让模型模仿人类执行任务的动作。但真实任务是\textbf{多步序列}，所以必须在多步轨迹上训练（这对应 System-1 的“怎么把动作连起来”）。

\textbf{统一动作空间（Unified Action Space，见表~\ref{tab:action}）}：把各平台语义等价的动作标准化，比如 Windows 上的 “click” 与手机上的 “tap” 统一，便于跨平台知识迁移；同时为每个平台保留专属可选动作。还定义两个\textbf{终止动作}：\texttt{Finished()} 表示任务完成，\texttt{CallUser()} 表示需要用户介入（如登录、验证）。

\textbf{动作轨迹采集}：两个来源——（1）\textbf{自标注数据}：自研标注工具捕获用户在 PC 上各类软件/网站的操作，先拟定任务指令、标注者审校后执行、再严格质量过滤；（2）\textbf{开源数据}：整合 MM-Mind2Web、GUIAct、AITW、AITZ、AndroidControl、GUI-Odyssey、AMEX 等，统一成同一动作空间格式。统计见表~\ref{tab:data}。

\textbf{提升定位能力}：定位数据比多步数据\textbf{更容易规模化}（主要靠元素的视觉与位置属性，可合成或抽取）。训练模型\textbf{直接预测要交互元素的坐标}：对每个有包围框的元素，取四角平均得到中心点；输出\textbf{归一化到屏幕尺寸的相对坐标}（保证不同分辨率设备一致）。例如给描述“右上角写着 Submit 的红按钮”，模型预测该按钮的归一化坐标。还整合了 SeeClick、MultiUI、Rico-SCA、CLAY、OmniACT、OS-ATLAS 等大量开源定位数据。

\begin{table}[H]\centering
\caption{统一动作空间（节选，对应原文 Table 1）。Shared 为跨平台共享动作，Desktop / Mobile 为各平台专属。}
\label{tab:action}
\small
\begin{tabular}{lll}
\toprule
平台 & 动作 & 含义 \\
\midrule
\multirow{7}{*}{Shared（共享）}
 & \texttt{Click(x, y)} & 在坐标 (x, y) 处点击 \\
 & \texttt{Drag(x1, y1, x2, y2)} & 从 (x1, y1) 拖拽到 (x2, y2) \\
 & \texttt{Scroll(x, y, direction)} & 在 (x, y) 处按方向滚动 \\
 & \texttt{Type(content)} & 输入指定内容 \\
 & \texttt{Wait()} & 短暂等待 \\
 & \texttt{Finished()} & 标记任务完成（终止） \\
 & \texttt{CallUser()} & 请求用户介入（终止） \\
\midrule
Desktop（桌面）
 & \texttt{Hotkey(key)} & 按下指定快捷键 \\
 & \texttt{LeftDouble(x, y)} & 在 (x, y) 双击 \\
 & \texttt{RightSingle(x, y)} & 在 (x, y) 右键单击 \\
\midrule
\multirow{4}{*}{Mobile（移动）}
 & \texttt{LongPress(x, y)} & 在 (x, y) 长按 \\
 & \texttt{PressBack()} & 按“返回”键 \\
 & \texttt{PressHome()} & 按“主页”键 \\
 & \texttt{PressEnter()} & 按“回车”键 \\
\bottomrule
\end{tabular}
\end{table}

\begin{table}[H]\centering
\caption{定位数据与多步动作轨迹数据的基本统计（对应原文 Table 2），对比作者自标注数据与开源数据。Ele.=元素数，Trace=动作轨迹数，avg steps=平均步数。“*”表示作者未披露具体数值。}
\label{tab:data}
\small
\begin{tabular}{llcccc}
\toprule
数据类型 & 平台 & Ele.（元素） & Ele./Image & Trace（轨迹） & avg steps \\
\midrule
\multirow{3}{*}{Open Source（开源）}
 & Web    & 14.8M & 6.7 & 6.4k & 7.1 \\
 & Mobile & 2.5M  & 4.6 & 145k & 9.6 \\
 & Desktop& 1.1M  & 6.2 & 0    & 0 \\
\midrule
Ours（本文标注） & --- & * & 7.5 & * & 14.9 \\
\bottomrule
\end{tabular}
\end{table}

\subsection{注入 System-2 推理（§4.4）}

只靠 System-1 直觉无法应对复杂多变的场景，因此要让模型学会\textbf{通过理解任务的全局结构来灵活规划动作步骤}。分两步走：

\textbf{第一步：用 GUI 教程做推理富化（Reasoning Enrichment）。} 网上图文交织的教程天然演示了各种软件/网页的详细交互，是建立 GUI 知识、引入逻辑推理模式的理想来源。作者从 MINT 和 OmniCorpus 两个图文交织数据集出发，做\textbf{多阶段过滤}：（1）\textbf{粗过滤}——训练一个 fastText 分类器筛出“像教程”的内容；（2）\textbf{细过滤}——用 LLM 剔除误报；（3）\textbf{去重与精炼}——用 URL 和 LSH（局部敏感哈希）去重，再让 LLM 改写文本。最终得到约 \textbf{600 万}条高质量 GUI 教程，平均每条 510 个文本 token、3.3 张图。

\textbf{第二步：用“思考增补”做推理激发（Thought Augmentation）。} §4.3 采集的轨迹只有“观测-动作”、\textbf{缺少显式思考}。于是给数据增补思考 $t$，把格式从 $(o_{i-1},a_{i-1},o_i,a_i,\dots)$ 变成 $(o_{i-1},t_{i-1},a_{i-1},o_i,t_i,a_i,\dots)$。两种标注方式：

\emph{(1) ActRe（反向标注）}：对每个轨迹按步拆分，对第 $n$ 步，用 VLM 在“已知前文上下文 + 当前目标动作 $a_n$”的条件下\textbf{反推}出思考 $t_n$：
\[
\begin{aligned}
t_n &= \mathrm{VLM}\big(\text{instruction},(o_1,t_1,a_1),(o_2,t_2,a_2),\dots,o_n,a_n\big)\\
t_{n+1} &= \mathrm{VLM}\big(\text{instruction},(o_1,t_1,a_1),\dots,(o_n,t_n,a_n),o_{n+1},a_{n+1}\big)
\end{aligned}
\]
这样生成的思考在逻辑上扎根于前文、且与当前动作对齐。标注时引导 VLM 表现出 System-2 的五种推理模式（对应原文 Figure 6）：\textbf{任务分解、长期一致性、里程碑识别、试错、反思}。

\emph{(2) 思考自举（Thought Bootstrapping，正向标注）}：上面的反向标注有个隐患——\textbf{先知道答案再编理由}，生成的思考可能只是“表面上”对得上动作，并非真正的因果推理（假阳性）。为此改用\textbf{不预先知道正确动作}的自举方式：让早期检查点模型 $\text{UI-TARS}_{\text{early}}$ 采样多个“思考-动作”对，再\textbf{挑出那个能导出正确动作的思考}：
\[
\begin{aligned}
(\hat t_n^{\,i},\hat a_n^{\,i})_{i=1}^{\text{max-try}} &= \text{UI-TARS}_{\text{early}}\big(\text{instruction},(o_1,t_1,a_1),\dots,o_n\big)\\
&\text{Select}(\hat t_n^{\,i},\hat a_n^{\,i}),\ \text{使得}\ \hat a_n^{\,i}=a_n
\end{aligned}
\]
这迫使模型\textbf{真正模拟决策过程}而非给既定动作找借口，标注质量更高。思考用\textbf{中英双语}标注以增加语言多样性；同时训练里也保留不带思考的原始轨迹。

\subsection{从长期记忆中的先验经验学习（§4.5）}

GUI 智能体难以像 LLM 那样规模化，根因是\textbf{大规模、标准化、真实世界的 GUI 操作过程数据极其稀缺}。解法：从长期记忆里的先验经验学习——让 UI-TARS \textbf{在真实设备上动态学习}，通过半自动的采集-过滤-精炼持续自我提升。

\textbf{在线轨迹自举（Online Trace Bootstrapping，对应原文 Figure 7）}：先拿到多样的任务目标（人工标注 + 模型生成）。在第 $n$ 轮，智能体 $M_n$ 在目标 GUI 环境（如虚拟 PC）里执行指令集 $I_n$，产生原始轨迹集 $T_{\text{raw},n}$。再用\textbf{多级过滤}得到干净轨迹：
\[
\mathrm{Filter}(T_{\text{raw},n},I_n)=T_{\text{filtered},n},
\]
过滤含三步：（1）\textbf{规则奖励}——启发式规则去掉明显异常（如不改变环境的冗余动作）；（2）\textbf{VLM 打分}——低于阈值的轨迹剔除；（3）\textbf{人工复核}——标注者找出出错步、丢弃其后动作、只保留有效前缀。然后用过滤后的轨迹自我提升、并扩充指令集：
\[
M_{n+1}=\mathrm{FineTune}(M_n,T_{\text{filtered},n}),\qquad
I_{n+1}=\mathrm{HumanRefine}(I_n,T_{\text{filtered},n}).
\]
在数百台虚拟 PC 上多轮迭代，不断用最新模型生成新轨迹。

\textbf{反思微调（Reflection Tuning）}：真实在线部署中，智能体常因不会自我反思而\textbf{卡死}（如反复点一个没反应的按钮）。而离线数据大多是“理想化、零错误”的轨迹，反而\textbf{教不会模型从错误中恢复}。为此提出反思微调：把模型自己犯的真实错误及其修正暴露给它学习。对一条出错的在线轨迹，设第 $\tau$ 步动作 $a_\tau$ 无效或次优，标注者给出修正后的思考与动作 $t^*_\tau, a^*_\tau$，得到\textbf{错误纠正对}：
\[
\begin{aligned}
T_- &= \big(\text{instruction},(o_1,t_1,a_1),\dots,(o_\tau,t_\tau,a_\tau)\big)\\
T_+ &= \big(\text{instruction},(o_1,t_1,a_1),\dots,(o_\tau,t^*_\tau,a^*_\tau)\big)
\end{aligned}
\]
更巧妙的是，还要求标注者\textbf{基于错误动作 $a_\tau$ 继续标注下一步}（模拟“错误已经发生”的情形）：下一步的思考 $t^*_{\tau+1}$ 必须\textbf{承认上一步失误、弥补其影响}，再给出能让任务重回正轨的 $a^*_{\tau+1}$（例如：上一步想加书签却误点了关闭，下一步就该重新打开刚关掉的网页再去点书签）。这构成\textbf{反思后（post-reflection）对}。训练时只用正样本 $T_+$ 做 SFT，且\textbf{只对修正步计损失}，错误步不参与训练。

\textbf{Agent DPO}：在线自举会自然产生大量\textbf{错误步骤（负样本）}，但 SFT 只用了修正的正样本、浪费了负样本。于是引入 \textbf{DPO}，同时利用修正动作与错误动作。设状态 $s_\tau$ 包含指令与到当前步的交互历史，原本错误动作 $a_\tau$ 被修正为偏好动作 $a'_\tau$。基于 Bradley-Terry 模型，成对偏好似然为：
\[
P_\theta(a'_\tau \succ a_\tau \mid s_\tau)=\frac{\exp\big(r_\theta(s_\tau,a'_\tau)\big)}{\exp\big(r_\theta(s_\tau,a_\tau)\big)+\exp\big(r_\theta(s_\tau,a'_\tau)\big)},
\]
其中 $r_\theta(s,a)$ 是学到的奖励函数，$a'_\tau \succ a_\tau$ 表示 $a'_\tau$ 比 $a_\tau$ 更受偏好。沿用 DPO 的推导（带 KL 约束的 RL 目标下的解析最优策略），直接在偏好数据上优化：
\[
\mathcal{L}_{\text{DPO}}(\theta)=-\,\mathbb{E}_\tau\left[\log\sigma\!\left(\beta\log\frac{\pi_\theta(a'_\tau|s_\tau)}{\pi_{\text{SFT}}(a'_\tau|s_\tau)}-\beta\log\frac{\pi_\theta(a_\tau|s_\tau)}{\pi_{\text{SFT}}(a_\tau|s_\tau)}\right)\right].
\]
\textbf{公式拆解：}$\pi_\theta$ 是待优化的智能体，$\pi_{\text{SFT}}$ 是 SFT 后的参考智能体，$\beta$ 控制二者偏离程度，$\sigma$ 是 sigmoid。直觉上：\textbf{提高修正动作 $a'_\tau$ 的概率、压低错误动作 $a_\tau$ 的概率}，从而隐式地用一个奖励函数把模型推向更好的决策。

\subsection{训练（§4.6）：三阶段流程}

为与 Aguvis、OS-Atlas 公平对比，沿用同样的 VLM 底座 \textbf{Qwen-2-VL}，总数据量约 \textbf{50B token}，分三阶段、逐步引入更高质量数据：
\begin{enumerate}[leftmargin=1.6em]
\item \textbf{继续预训练（Continual Pre-training）}：用 §4 的全部数据（不含反思微调数据），恒定学习率，打牢感知/定位/动作轨迹的基础。
\item \textbf{退火（Annealing）}：选取感知、定位、动作轨迹、反思微调数据的\textbf{高质量子集}做退火，逐步收敛学习动态，得到 \textbf{UI-TARS-SFT}。
\item \textbf{DPO 阶段}：用在线自举得到的反思配对数据做 DPO，强化最优动作、惩罚次优动作，得到最终的 \textbf{UI-TARS-DPO}。
\end{enumerate}

% =========================================================
\section{方法图解：UI-TARS 整体架构（图 4）}

\begin{figure}[H]\centering
\includegraphics[width=0.92\linewidth]{cand_p14_k00.pdf}
\caption{UI-TARS 总体架构（对应原文 Figure 4）。图中展示了模型架构及其四大核心能力。}
\label{fig:arch}
\end{figure}

\textbf{这张图在说什么、数据怎么流动？}整张图分三层：

\begin{itemize}[leftmargin=1.5em]
\item \textbf{顶层 Environment（环境）}：左边是输入——\textbf{User Query（用户指令）}和\textbf{动作空间}。中间一台电脑表示 GUI 环境。数据流是一个闭环：模型输出动作经 \textbf{PyAutoGUI}（一个能模拟真实鼠标键盘的库）作用到环境上 $\rightarrow$ 环境返回新的 \textbf{Observation（截图观测）}。沿时间轴排开的一串小卡片就是 $(o_i,t_i,a_i)$ 序列：每张卡片含一帧截图、一段 \textbf{Thought（思考，图中红字）}和一个 \textbf{Action（动作）}——比如“我意识到打开了 Word，需要双击桌面图标……”然后执行 \texttt{left\_double(...)}。这正对应 §4.1 的两条公式：动作前先想清楚。
\item \textbf{中层 UI-TARS}：就是\textbf{那一个端到端模型}本身，承接环境的输入、吐出思考与动作。强调“原生”——不靠外挂模块。
\item \textbf{底层 Capability（能力）}：把全文方法的四块对应起来——\textbf{Perception（感知）}列出五类任务（元素描述/密集描述/转移描述/问答/SoM，§4.2）；\textbf{Action（动作）}展示统一动作空间 + 自标注数据 + 开源数据 $\rightarrow$ 多步轨迹数据（§4.3）；\textbf{System-2 Reasoning}对应教程富化 + 思考增补（§4.4）；\textbf{Learning from Prior Experience（从先验经验学习）}对应在线轨迹自举 + 反思微调 + Agent DPO（§4.5）。
\end{itemize}

一句话：\textbf{上面是“怎么和世界交互”（看截图→想→动→再看），下面是“靠哪四种能力把这件事做好”，中间是把这一切融成一体的那个模型。}

% =========================================================
\section{实验解读（§5）}

\textbf{模型与基线。} 三个规模：UI-TARS-2B / 7B / 72B（底座 Qwen-2-VL）。基线覆盖\textbf{商业模型}（GPT-4o、Claude-3.5-Sonnet、Gemini-1.5-Pro、Gemini-2.0 即 Project Mariner、Claude Computer-Use）与\textbf{学术模型}（CogAgent、OmniParser、InternVL、Aria-UI、Aguvis、OS-Atlas、UGround、ShowUI、SeeClick、Qwen 系列等）。下文挑关键数字重排成小表。

\subsection{感知能力（§5.1）}

三个基准：VisualWebBench（网页元素理解与定位）、WebSRC（网页结构理解 QA）、ScreenQA-short（移动屏幕内容理解 QA）。

\begin{table}[H]\centering
\caption{GUI 感知基准结果（节选自原文 Table 3，数值越高越好；\textbf{粗体}为该列最佳）。}
\label{tab:perc}
\small
\begin{tabular}{lccc}
\toprule
模型 & VisualWebBench & WebSRC & ScreenQA-short \\
\midrule
Claude-3.5-Sonnet & 78.2 & 90.4 & 83.1 \\
GPT-4o            & 78.5 & 87.7 & 82.3 \\
UI-TARS-2B        & 72.9 & 89.2 & 86.4 \\
UI-TARS-7B        & 79.7 & \textbf{93.6} & 87.7 \\
UI-TARS-72B       & \textbf{82.8} & 89.3 & \textbf{88.6} \\
\bottomrule
\end{tabular}
\end{table}

\textbf{说明什么：}UI-TARS-72B 在 VisualWebBench 上 82.8 全面超过 GPT-4o（78.5）和 Claude（78.2）；7B 在 WebSRC 上拿到 93.6 领先。\textbf{看界面看得准，是后面能点对、做对任务的地基。}

\subsection{定位能力（§5.2）}

三个基准：ScreenSpot Pro（高分辨率专业桌面环境，23 个应用、5 个行业、3 种操作系统，\textbf{最难}）、ScreenSpot、ScreenSpot v2。

\begin{table}[H]\centering
\caption{定位能力关键结果（综合自原文 Table 4--6，报告平均分 Avg）。}
\label{tab:ground}
\small
\begin{tabular}{lccc}
\toprule
模型 & ScreenSpot Pro & ScreenSpot & ScreenSpot v2 \\
\midrule
OS-Atlas-7B        & 18.9 & 82.5 & 87.1 \\
UGround-V1-7B      & 31.1 & ---  & ---  \\
Claude Computer Use& 17.1 & 83.0 & ---  \\
Aguvis-72B         & ---  & 89.2 & ---  \\
UI-TARS-2B         & 27.7 & 82.3 & 84.7 \\
UI-TARS-7B         & 35.7 & \textbf{89.5} & \textbf{91.6} \\
UI-TARS-72B        & \textbf{38.1} & 88.4 & 90.3 \\
\bottomrule
\end{tabular}
\end{table}

\textbf{说明什么：}在最难的 ScreenSpot Pro 上，UI-TARS-72B 取得 38.1 的 SOTA，远超 UGround-V1-7B（31.1）和 OS-Atlas-7B（18.9）。一个有趣发现：\textbf{提高输入图像分辨率能在 ScreenSpot Pro 上显著涨点}；而 ScreenSpot v1/v2 在 7B$\rightarrow$72B 几乎不涨，说明这两个老基准可能\textbf{不够难、难以充分体现大模型的定位上限}。

\subsection{离线智能体能力（§5.3）}

三个基准：Multimodal Mind2Web（网页，指标 Ele.Acc/Op.F1/Step SR）、AndroidControl（移动，分 Low/High 两档）、GUI Odyssey（移动跨应用，平均 15+ 步）。

\textbf{关键结论：}（1）在 Mind2Web 上，\textbf{智能体模型普遍碾压框架方法}（用 GPT-4o/GPT-4V 当规划器的）；UI-TARS-72B 多项指标 SOTA，连 7B 都超过了 Aguvis-72B 和 Claude。（2）在 AndroidControl 和 GUI Odyssey 上，UI-TARS-7B/72B 比前 SOTA（OS-Atlas-7B）\textbf{绝对提升约 25 个点}。（3）\textbf{Claude Computer-Use 在网页强、在手机上明显吃力}，说明其 GUI 能力没很好迁移到移动端；而 UI-TARS 网页和移动\textbf{两端都强}，体现更好的适应性与泛化。

\subsection{在线智能体能力（§5.4）}

两个动态交互基准：\textbf{OSWorld}（Ubuntu/Windows/macOS 桌面，369 个真实任务，仅截图模式，3 次平均）、\textbf{AndroidWorld}（真实 Android 模拟器，116 个任务、20 个应用）。

\begin{table}[H]\centering
\caption{在线基准结果（节选自原文 Table 9）。OSWorld 为仅截图设置，括号内为步数预算。}
\label{tab:online}
\small
\begin{tabular}{lcc}
\toprule
方法 & OSWorld & AndroidWorld \\
\midrule
GPT-4o（Agent Model） & 5.0 & 34.5 (SoM) \\
GPT-4o + Aria-UI（框架） & 15.2 & 44.8 \\
Aguvis-72B & 10.3 & 26.1 \\
Claude Computer-Use & 14.9 (15 步) / 22.0 (50 步) & 27.9 \\
UI-TARS-7B-DPO  & 18.7 (15 步) & 33.0 \\
UI-TARS-72B-SFT & 18.8 (15 步) & \textbf{46.6} \\
UI-TARS-72B-DPO & \textbf{22.7 (15 步) / 24.6 (50 步)} & --- \\
\bottomrule
\end{tabular}
\end{table}

\textbf{说明什么：}（1）OSWorld 15 步预算下，UI-TARS-72B-DPO（22.7）大幅超过 Claude（14.9）；而且\textbf{UI-TARS 用 15 步（22.7）就追平了 Claude 用 50 步（22.0）}——执行效率高。50 步下 UI-TARS-72B-DPO 拿到 24.6 的新 SOTA，超过所有框架方法。（2）AndroidWorld 上 UI-TARS-72B-SFT 得 46.6，超过最强框架（GPT-4o+Aria-UI，44.8）和最强模型（Aguvis-72B，26.1）。（3）\textbf{DPO 显著提升 OSWorld}，证明“训练时引入负样本”能让模型更好区分优/次优动作。（4）\textbf{72B 在在线任务上比 7B 强很多，且差距比离线大}——说明\textbf{放大模型规模能显著提升 System-2 推理}；也说明\textbf{光看离线基准会低估模型在真实动态环境里的能力}。

\subsection{System-1 vs.\ System-2 对比（§5.5）}

同一个 UI-TARS-7B 同时学会两种模式，推理时用提示工程切换。用 \textbf{Best-of-N}（$N=1,16,64$）采样，以 step 成功率衡量（对应原文 Figure 8）。

\textbf{两个反直觉但重要的发现：}
\begin{itemize}[leftmargin=1.5em]
\item \textbf{域内（in-domain）}：当 $N=1$（只采一次）时，System-2 反而\textbf{略逊于} System-1——因为多步推理在单样本下可能引入无关或错误的推理步（甚至指向不存在的对象），增加幻觉风险，模型容易“一条道走到黑”。但当 $N$ 增到 16、64，\textbf{System-2 明显反超}：候选多了，就有机会跳出糟糕的推理路径，正确动作落在采样里的概率大增。
\item \textbf{域外（OOD，AndroidWorld，无对应训练数据）}：即便在 Bo1（$N=1$）下，\textbf{System-2 也显著超过 System-1}。说明更深的推理在\textbf{没见过的新任务}上泛化优势明显。
\end{itemize}

\textbf{核心洞见：}System-2 在“多样性充足”时很强，但要在\textbf{单次、果断输出}下也达到最优仍是挑战。理想方向是用\textbf{强化微调（reinforced fine-tuning）}让模型一次就高置信地给出正确动作——这也为后续 UI-TARS-1.5/2 的强化学习路线埋下伏笔。

% =========================================================
\section{核心 Takeaways 与对 AI 应用开发的启示}

\begin{enumerate}[leftmargin=1.6em]
\item \textbf{“一个端到端模型”打败“强模型+人手工拼的框架”。}UI-TARS 证明：与其拿 GPT-4o 当大脑再外挂一堆提示词和工具模块，不如把感知/推理/记忆/动作\textbf{全部训进一个模型}。框架是设计驱动、脆弱难维护；原生模型是数据驱动、可持续进化。\emph{启示：}做 GUI/桌面自动化产品时，长期看\textbf{数据驱动的专用模型}比堆叠提示工程更可扩展。

\item \textbf{纯视觉 + 统一动作空间 = 跨平台通吃。}只看截图、把 click/tap 等统一抽象、输出归一化坐标，让同一套能力在网页、桌面、手机间迁移。\emph{启示：}设计自动化 agent 时，\textbf{统一的动作抽象层}能极大降低多端适配成本。

\item \textbf{“先想后做”（System-2 thought）确实有用，但要配合采样多样性。}动作前显式生成思考能提升复杂/陌生任务表现；但单次采样时慢思考可能帮倒忙，需要 Best-of-N 或强化微调来兑现其潜力。\emph{启示：}在产品里加 CoT 不是免费午餐，要权衡\textbf{延迟、采样次数与收益}。

\item \textbf{把“失败轨迹”当宝：反思微调 + Agent DPO。}离线理想数据教不会模型纠错；让模型在真实环境里犯错、标注修正、甚至\textbf{标注“错了之后怎么补救”}，再用 DPO 同时利用正负样本。\emph{启示：}\textbf{错误数据和成功数据一样有价值}，尤其对需要在线鲁棒性的 agent。

\item \textbf{在线自举突破数据瓶颈。}在数百台虚拟机上自动跑、规则+VLM+人工三级过滤、回灌迭代，是解决“GUI 过程数据稀缺”的工程范式。\emph{启示：}\textbf{半自动数据飞轮}（采集→过滤→训练→再采集）是 agent 类产品冷启动后持续提升的可复制套路。

\item \textbf{规模化对“慢思考”的回报更大；离线基准会低估真实能力。}72B 在在线动态任务上对 7B 的优势远大于离线，且要在真实交互环境里才看得出来。\emph{启示：}评估 agent \textbf{一定要上在线/真实环境}，别只看离线单步指标；推理密集型场景值得投更大模型。
\end{enumerate}

\section*{参考文献（精选）}
\begin{enumerate}[leftmargin=2em]
\item Yujia Qin, Yining Ye, Junjie Fang, et al. UI-TARS: Pioneering Automated GUI Interaction with Native Agents. arXiv preprint arXiv:2501.12326, 2025. \url{https://arxiv.org/abs/2501.12326}
\item Peng Wang, Shuai Bai, Sinan Tan, et al. Qwen2-VL: Enhancing Vision-Language Model's Perception of the World at Any Resolution. 2024.
\item Rafael Rafailov, Archit Sharma, Eric Mitchell, et al. Direct Preference Optimization: Your Language Model is Secretly a Reward Model. NeurIPS, 2023.
\item Shunyu Yao, Jeffrey Zhao, Dian Yu, et al. ReAct: Synergizing Reasoning and Acting in Language Models. ICLR, 2023.
\item Tianbao Xie, Danyang Zhang, Jixuan Chen, et al. OSWorld: Benchmarking Multimodal Agents for Open-Ended Tasks in Real Computer Environments. NeurIPS, 2024.
\item Christopher Rawles, Sarah Clinckemaillie, Yifan Chang, et al. AndroidWorld: A Dynamic Benchmarking Environment for Autonomous Agents. 2024.
\item Anthropic. Developing a Computer Use Model. 2024. \url{https://www.anthropic.com/news/developing-computer-use}
\end{enumerate}

\end{document}
